% =============================================================
% §7 Related Work
% =============================================================
\section{Related Work}
\label{sec:related}

\textbf{LLM serving systems (single server or cluster).}
Orca~\cite{yu2022orca} introduced iteration-level continuous
batching, vLLM~\cite{kwon2023vllm} added PagedAttention
for non-contiguous KV allocation
(2--4$\times$ throughput), and
Sarathi-Serve~\cite{agrawal2024sarathi} introduced chunked
prefill to mitigate generation stalls in mixed batches.
VoltanaLLM~\cite{yu2025voltana} co-optimizes energy and \slo{}
inside a prefill/decode-disaggregated cluster via frequency
control and state-space routing. These systems optimize
\emph{within} one serving engine or cluster and are
\textbf{orthogonal} to our cross-server scheduling: our
admission-time batching is a tractable abstraction of their
iteration-level batching, and our per-request placement
decisions would propagate into vLLM-style engines unchanged.

\textbf{Cluster scheduling: from heuristics to learning.}
Classic schedulers---HEFT~\cite{topcuoglu2002heft} with
upward-rank priorities, and metaheuristics such as
GA~\cite{holland1975adaptation} and
PSO~\cite{kennedy1995pso}---optimize fixed cost models for
generic DAG workloads; industry favors simpler load-balancing
primitives. Learned schedulers followed:
Decima~\cite{mao2019decima} (GNN + REINFORCE for Spark DAGs),
RLScheduler~\cite{zhang2020rlscheduler} (HPC batch jobs), and
A3C-R2N2~\cite{tuli2022a3cr2n} (actor-critic for generic
edge-cloud tasks, predating the LLM era)---the closest
predecessor in deployment topology, which we benchmark as a
\textbf{controlled baseline} (\S\ref{sec:main-results}) to
isolate our AIGC contribution. This line remains active:
recent work applies GNN-based DRL to energy-aware cloud
scheduling~\cite{dagtopology2026} and diffusion-aided reward
shaping to AIGC workload scheduling across data
centers~\cite{aigcenergy2026}. Two findings from this
literature prefigure ours: Decima's gains concentrate at high
cluster load, and \cite{dagtopology2026} observes that each
learned policy degenerates to a regime-specific
heuristic---consistent with our regime characterization (C4).
None of these schedulers encodes LLM-serving physics.

\textbf{LLM request scheduling across servers.} A fast-growing
line schedules LLM inference \emph{across} instances rather
than within one engine. Inside a single cluster, Microsoft's
workload-aware router~\cite{jain2024router} trains a
heuristic-guided RL policy over prefill/decode and batch-mixing
effects, and Lodestar~\cite{lim2026lodestar} learns an online
reward predictor coupling batching with KV-cache reuse---both
echo our thesis that serving physics should be exposed to a
learned router, but neither models a cloud-edge hierarchy nor
energy. FREESH~\cite{he2025freesh} routes and schedules for
energy and carbon across geo-distributed heterogeneous GPUs
through classical optimization, and
MARLIN~\cite{moore2026marlin} applies multi-agent
game-theoretic RL to multi-objective LLM scheduling within
cloud datacenters. On the cloud-edge side,
PerLLM~\cite{yang2024perllm} is the closest prior work: it
schedules diverse LLM service requests across edge and cloud
with a constrained combinatorial bandit and reports large
throughput and energy-cost gains. We differ on three axes:
\emph{(i)~method}---PerLLM's bandit optimizes a scalarized
cost under constraints, whereas we learn a deep policy whose
\textbf{state and reward expose serving physics} (model
residency, KV-cache locality, batch occupancy);
\emph{(ii)~objective}---we treat \slo{} and energy per token
as an explicit Pareto pair rather than folding them into one
cost; \emph{(iii)~analysis}---our joint ablation and
weight-sensitivity study (\S\ref{sec:ablation},
\S\ref{sec:weight-sens}) quantify \emph{where} such gains come
from, an analysis bandit-style methods do not expose. We
reimplement a CS-UCB-style scheduler in our simulator as a
baseline; it clusters with the generic baselines on the
Pareto plane (\S\ref{sec:main-results}). Other
cloud-edge efforts differ in decision granularity or method:
Lyapunov-assisted DRL partitions transformer layers across
tiers~\cite{younesi2025lyapunov}; active-inference
offloading~\cite{he2024activeinference} and
world-model-assisted PPO offloading with on-device
compression~\cite{zhang2026worldmodel} make device-side
offload decisions; T2DRL~\cite{xu2025t2drl} co-designs model
caching and offloading at the mobile edge; QoS-aware routing
selects among edge-deployed experts~\cite{yang2025qosrouting};
NSGA-II search routes requests across cloud-edge LLM instances
without learning~\cite{yu2025llmrouting}; and
SLICE~\cite{chow2025slice} schedules differentiated \slo{}
tiers on a single edge device via utility heuristics and
decode-rate control. To our knowledge, \textbf{no prior work
combines a deep-RL policy over a heterogeneous cloud-edge
pool with serving-physics state/reward and joint
\slo{}+energy Pareto evaluation}---the specific conjunction
this paper studies and ablates.

\textbf{AIGC deployment topologies.}
Splitwise~\cite{patel2024splitwise} and
DistServe~\cite{zhong2024distserve} disaggregate prefill from
decode across machines of different generations;
Llumnix~\cite{sun2024llumnix} migrates requests across
instances mid-flight. These works design \emph{deployment
topologies} and \emph{reactive policies}, treating initial
request placement as orthogonal---the two strands compose:
our scheduler could populate prefill/decode pools in a
Splitwise-style deployment, and Llumnix-style migration could
refine our placements at runtime.

\textbf{Simulation platforms.}
CloudSim~\cite{calheiros2011cloudsim},
iFogSim~\cite{gupta2017ifogsim}, and
WorkflowSim~\cite{chen2012workflowsim} target general
cloud-fog or DAG workloads without LLM physics. A recent wave
of LLM-serving simulators models intra-cluster serving with
high fidelity: Vidur~\cite{agrawal2024vidur} builds
profiling-plus-ML operator models for cluster-scale what-if
analysis; SplitwiseSim~\cite{patel2024splitwise} simulates
phase-split clusters; LLMServingSim~\cite{cho2024llmservingsim}
co-simulates hardware and software;
Frontier~\cite{frontier2025sim} models disaggregated serving
with explicit cross-server KV-transfer events; and Helix's
simulator covers geo-distributed heterogeneous GPU
pools~\cite{mei2025helix}. An extension of Vidur couples a GPU
power model to quantify energy and
carbon~\cite{wiesner2025vidurvessim}. Each pillar of our
simulator therefore exists somewhere in isolation---continuous
batching and the two-phase lifecycle are standard; Frontier
models cross-server KV movement;
\cite{wiesner2025vidurvessim} adds energy accounting; and
weight-residency/cold-load dynamics appear in serverless
serving systems~\cite{fu2024serverlessllm}. To our knowledge,
however, ours is the first \emph{open} platform to
\textbf{combine} all five AIGC physics with a cloud-edge
cross-server scheduling scope and native multi-objective
(\slo{}~+~energy) evaluation. We release it with
reproducibility manifests for the 30+ configurations in
\S\ref{sec:eval}.
